\documentclass[a4paper, 11pt]{article}

% Language and encoding.
\usepackage[english]{babel}
\usepackage[utf8]{inputenc}

% Set fonts (in order to deal with umlauts).
\usepackage[T1]{fontenc}
\usepackage{lmodern}

% Sets page size and margins.
\usepackage[a4paper, top=2.5cm, bottom=2.5cm, left=3cm, right=3cm, marginparwidth=1.75cm]{geometry}

% Useful packages
\usepackage{graphicx}
\usepackage{subcaption}
\usepackage{url}

% Creative Commons license package.
\usepackage[
    type={CC},
    modifier={by-nc-sa},
    version={4.0},
]{doclicense}

% hyperref package to color links.
%\usepackage[colorlinks=true, allcolors=blue]{hyperref}

% Todo notes.
%\usepackage[textsize=footnotesize, backgroundcolor=yellow!10, linecolor=gray!35]{todonotes} 

\title{Gender Differences and Bias in Open Source \& Open Data}

\begin{document}

\date{\today}
\author{E. Jobst and M. Zandpour}
\maketitle

\begin{abstract}
Some abstract...
\end{abstract}

\section{Introduction} \label{sec:intro}

\subsection{Motivation} \label{sec:intro:motivation}

\subsection{Summary} \label{sec:intro:summary}

\section{Gender Disparity in Computer Science} \label{sec:gender-disparity}

\subsection{Introduction} \label{sec:gender-disparity:introduction}

\subsection{Contributions to Lack of Female Participation} \label{sec:gender-disparity:contributions}

\subsection{Efforts in Increasing Female Participation} \label{sec:gender-disparity:efforts}

\subsection{Gender Disparity in Open Source \& Open Data}

\section{Case Studies} \label{sec:casestudies}

\subsection{Wikipedia} \label{sec:casestudies:wikipedia}

\subsection{GitHub} \label{sec:casestudies:github}

\subsection{3rd Case Study}

\subsection{4th Case Study}

\section{Conclusion} \label{sec:conclusion}

% References section.
\bibliographystyle{plain}
\bibliography{bibliography}

% Add CC license.
\doclicenseThis

\section{Latex QuickTips} \label{sec:tips}

To ensure best compatibility with version control (e.g.\ git), it is best to write one sentence per line and configure your editor not to hard-break lines.
For citations use \cite{terrell-gender-bias} or \cite[p.\ 6--8]{terrell-gender-bias}.
For quotes use, for example,
\begin{quote}
	``This is a proper quote.'' 
\end{quote}

\begin{figure}[t]
	\centering
	\includegraphics[width=0.25\textwidth]{figures/logo.png}
\caption{LaTeX logo.} \label{fig:logo}
\end{figure}

Sometimes, we would like to include figures.
This sentence provides an example of referencing to Fig.~\ref{fig:logo}.

You can of course also refer to Sec.~\ref{sec:intro} or to Sec.~\ref{sec:intro:motivation}.

Footnotes are includes, for example, here.\footnote{A footnote containing a link \url{https://free-and-open-technologies.github.io}}

\end{document}
